\documentclass{article}

\usepackage{amsmath,amssymb}
\usepackage{xurl}
\usepackage[hidelinks]{hyperref}
\setlength{\emergencystretch}{2em}

% Shared commands for notes in docs/paper-gaps/.

\DeclareUnicodeCharacter{2080}{\ifmmode{}_{0}\else\textsubscript{0}\fi}
\DeclareUnicodeCharacter{2081}{\ifmmode{}_{1}\else\textsubscript{1}\fi}
\DeclareUnicodeCharacter{2082}{\ifmmode{}_{2}\else\textsubscript{2}\fi}
\DeclareUnicodeCharacter{2099}{\ifmmode{}_{n}\else\textsubscript{n}\fi}

\newcommand{\Lean}{\textsc{Lean}}
\ExplSyntaxOn
\NewDocumentCommand{\leanid}{m}{
  \ifmmode
    \textsf{\detokenize{#1}}
  \else
    \group_begin:
    \urlstyle{sf}
    \tl_set:Nn \l_tmpa_tl {#1}
    \tl_replace_all:Nnn \l_tmpa_tl {\_} {_}
    \exp_args:NV \path \l_tmpa_tl
    \group_end:
  \fi
}
\ExplSyntaxOff
\providecommand{\tr}{\operatorname{tr}}
\newcommand{\tnleanIssueBase}{https://github.com/LionSR/TNLean/issues/}
\newcommand{\tnleanPrBase}{https://github.com/LionSR/TNLean/pull/}
\newcommand{\tnleanDocBase}{https://sirui-lu.com/TNLean/docs/}
\newcommand{\ghissue}[1]{\href{\tnleanIssueBase#1}{GitHub issue \##1}}
\newcommand{\ghpr}[1]{\href{\tnleanPrBase#1}{PR \##1}}
\newcommand{\doclink}[1]{\href{\tnleanDocBase#1.html}{\path{#1}}}

% Dirac notation and the identity operator, as in blueprint/src/macros/common.tex.
% \providecommand keeps note-local definitions authoritative.
\usepackage{dsfont}
\providecommand{\ket}[1]{|#1\rangle}
\providecommand{\bra}[1]{\langle #1|}
\providecommand{\braket}[2]{\langle #1 | #2 \rangle}
\providecommand{\ketbra}[2]{|#1\rangle\!\langle #2|}
\providecommand{\Id}{\mathds{1}}
\providecommand{\one}{\mathds{1}}
\providecommand{\C}{\mathbb{C}}
\providecommand{\MN}[1]{M_{#1}(\C)}
\providecommand{\MD}{M_D(\C)}

% External referencing for paper sources.  The build helper writes sanitized
% auxiliary files under build/paper-aux/ and excludes duplicated source labels.
\usepackage{xr}

% Import a generated paper auxiliary file with a caller-supplied label prefix.
% The candidate paths support builds from docs/paper-gaps/, the repository root,
% and build/paper-aux/ itself.
\newcommand{\importpaperaux}[2]{%
  \IfFileExists{../../build/paper-aux/#2.aux}{%
    \externaldocument[#1]{../../build/paper-aux/#2}%
  }{%
    \IfFileExists{build/paper-aux/#2.aux}{%
      \externaldocument[#1]{build/paper-aux/#2}%
    }{%
      \IfFileExists{#2.aux}{%
        \externaldocument[#1]{#2}%
      }{}%
    }%
  }%
}

% General external reference macros
\newcommand{\paperref}[2]{\ref{#1-#2}}
\newcommand{\papereqref}[2]{(\ref{#1-#2})}

% CPSV16 source (1606.00608 / MPDO-22-12-17-2.tex) external document and shortcuts
\importpaperaux{cpsv16-}{1606.00608/MPDO-22-12-17-2-tnlean}
\newcommand{\cpsvref}[1]{\paperref{cpsv16}{#1}}
\newcommand{\cpsveqref}[1]{\papereqref{cpsv16}{#1}}

% Verdict marker: \gapnote{<kind>}{<status>}. Kind is the policy
% classification (clarification, local-correction, scope-restriction,
% unfaithful, false-source, open-gap); status is open, wip (elimination
% underway), resolved, or historical. Machine-read by the site index;
% prints a verdict line.
\newcommand{\gapnote}[2]{\par\noindent\textbf{Verdict:} #1 (#2).\par}


\title{Deferred Result: The Translation-Invariant Bond-Dimension Lower Bound
for the W State}
\author{}
\date{2026-06-17, revised 2026-09-26}

\begin{document}
\maketitle
\gapnote{open-gap}{open}

\section{Source statement}

For $N\geq 1$, define the unnormalized $N$-site W state by
\begin{align}
    \ket{W_N}
        &=
        \ket{10\cdots 0}
        +\ket{010\cdots 0}
        +\cdots
        +\ket{0\cdots 01}.
    \label{eq:rmp_w_state_ti_bound_w_state}
\end{align}
The review presents $\ket{W_N}$ as an open-boundary matrix product state with
bond dimension $D=2$ \cite[lines~2352--2361]{Cirac2021Matrix}. Its
site-independent matrices are
\begin{align}
    A^0
        &=
        \begin{pmatrix}
            1 & 0\\
            0 & 1
        \end{pmatrix},
    \label{eq:rmp_w_state_ti_bound_zero_matrix}\\
    A^1
        &=
        \begin{pmatrix}
            0 & 1\\
            0 & 0
        \end{pmatrix}.
    \label{eq:rmp_w_state_ti_bound_one_matrix}
\end{align}
With boundary covectors $(l|=(0|$ and $|r)=|1)$, these matrices give the
open-boundary representation.

This representation uses a site-independent bulk tensor, but its boundary
contraction is not a translation-invariant periodic representation. The review
states that any translation-invariant matrix product state whose periodic
contraction equals $\ket{W_N}$ must have a bond dimension that grows with
$N$ \cite[line~2362]{Cirac2021Matrix}. More precisely, it states a bound of
the form
\begin{align}
    D^3\log D
        &=
        \Omega(N).
    \label{eq:rmp_w_state_ti_bound_asymptotic_lower_bound}
\end{align}
Consequently, for every $\delta>0$,
\begin{align}
    D
        &=
        \Omega\left(N^{1/(3+\delta)}\right).
    \label{eq:rmp_w_state_ti_bound_polynomial_consequence}
\end{align}

\section{Source proof}

The review does not prove
(\ref{eq:rmp_w_state_ti_bound_asymptotic_lower_bound}) directly. It combines
the corollary on the W state in Ref.~\cite[lines~2182--2225]{PerezGarcia2007MPSReps}
with the quantum Wielandt inequality \cite{Sanz2010Quantum} and its refinement
attributed there to Micha\l{}ek and Shitov \cite[line~2362]{Cirac2021Matrix}.

The corollary of Ref.~\cite{PerezGarcia2007MPSReps} states: if $\ket{W_N}$ is
the periodic vector of $D\times D$ matrices, and Conjecture~2 of that paper
holds, then $D\succeq O(N^{1/3})$. Conjecture~2
\cite[lines~2109--2111]{PerezGarcia2007MPSReps} asserts that condition C1
(the words of length $L_0$ span the full matrix algebra) holds in every block
of the canonical form with $L_0=O(D^2)$. The proof runs as follows.
\begin{enumerate}
    \item Let the canonical form have $b$ pairwise different blocks of sizes
    $D_1,\ldots,D_b$ and weights $\lambda_j$, so that
    $\ket{W_N}=\sum_j\lambda_j^N\sum_\sigma
    \tr(A^j_{\sigma_1}\cdots A^j_{\sigma_N})\ket{\sigma}$.
    Assume $N/2\geq 3(b-1)(L_0+1)$.

    \item Cut the chain into two pieces of lengths $R,R'$, each at least
    $3(b-1)(L_0+1)$. By the direct-sum lemma
    \cite[lines~1346--1349]{PerezGarcia2007MPSReps} and condition C1 in each
    block, the reduced state of the first piece has rank at least
    $\sum_jD_j^2$.

    \item The reduced state of $\ket{W_N}$ across any cut has rank $2$.
    Hence $\sum_jD_j^2\leq 2$, so there are at most two blocks, each of size
    $1\times1$.

    \item The source calls it trivial that this is impossible. Therefore
    $N/2\leq 3(b-1)(L_0+1)$, and with $L_0=O(D^2)$ and $b\leq D$ the bound
    $D\succeq O(N^{1/3})$ follows.
\end{enumerate}
The review replaces Conjecture~2 by the Wielandt bound $L_0=O(D^2\log D)$ of
Micha\l{}ek and Shitov, which gives
(\ref{eq:rmp_w_state_ti_bound_asymptotic_lower_bound}).

\section{Formalized conditional core}

Steps~1--4 are formalized for a tensor in the translation-invariant canonical
form whose blocks satisfy condition C1 at a common length $L_0>0$ and are
pairwise distinct in the sense used by the direct-sum lemma: no two blocks of
equal size are related by a gauge transformation and a phase. The formal
results are the following.
\begin{itemize}
    \item The amplitudes of a canonical form are the weighted sums of the block
    amplitudes (step~1).\footnote{Formal declarations:
    \leanid{MPSTensor.PGVWC07CanonicalFormData.trace_evalWord_eq_sum} and
    \leanid{MPSTensor.PGVWC07CanonicalFormData.mpv_eq_sum}.}

    \item For every length $n\geq\max(L_0,3(b-1)(L_0+1))$ the tuples of block
    words of length $n$ span the product of the block matrix algebras. Condition
    C1 in a block also makes the block irreducible.\footnote{Formal
    declarations:
    \leanid{MPSTensor.PGVWC07CanonicalFormData.wordTupleSpanTop_of_ge},
    \leanid{Kraus.isIrreducibleFamily_of_isNBlkInjective}, and
    \leanid{MPSTensor.wordTupleSpanTop_of_le_one}.}

    \item The rank inequality of step~2: if the functions of the first piece
    obtained by fixing a configuration of the second piece lie in a subspace
    $V$, then $\sum_jD_j^2\leq\dim V$.\footnote{Formal declaration:
    \leanid{MPSTensor.sum_sq_dim_le_finrank_of_forall_mem}.}

    \item The cut rank of $\ket{W_N}$ is at most $2$ (step~3): the amplitude
    of a concatenation $\sigma\tau$ is
    $u_1(\sigma)u_0(\tau)+u_0(\sigma)u_1(\tau)$, where $u_m$ indicates
    configurations with $m$ excitations.\footnote{Formal declarations:
    \leanid{MPSTensor.wIndicator_append},
    \leanid{MPSTensor.wIndicator_append_mem_span}, and
    \leanid{MPSTensor.finrank_span_wIndicator_append_le_two}.}

    \item The step called trivial (step~4). Let $x_k,y_k$ be the weighted
    one-site amplitudes of at most two $1\times1$ blocks, so that the
    amplitude of $1^m0^{N-m}$ is $g_m=\sum_ky_k^mx_k^{N-m}$. If
    $g_1=c\neq0$ and $g_2=g_3=0$, the identity
    $x_1x_2g_3-(x_1y_2+x_2y_1)g_2+y_1y_2g_1=0$ gives $y_1y_2=0$. Then $g_1$ and $g_2$ reduce to the terms of the
    remaining block $k$, and $c\,y_k=y_k^2x_k^{N-1}=x_kg_2=0$, so $y_k=0$ and
    $g_1=0$, a contradiction.
    This needs $N\geq3$ for two blocks and $N\geq2$ for one
    block.\footnote{Formal declaration:
    \leanid{MPSTensor.not_sum_pow_mul_pow_eq}.}

    \item The conclusion:\footnote{Formal declaration:
    \leanid{MPSTensor.PGVWC07CanonicalFormData.lt_of_mpv_eq_smul_wIndicator}.}
    if the periodic vector on $N$ sites is $c\ket{W_N}$ with $c\neq0$, then
    \begin{align}
        N
            &<
            2\max\bigl(L_0,\,3(b-1)(L_0+1)\bigr).
        \label{eq:rmp_w_state_ti_bound_conditional}
    \end{align}

    \item An explicit bound in $D$, weaker than the source's. If the adjoint
    channel $X\mapsto\sum_iA^{j\dagger}_iXA^j_i$ of every block is primitive,
    Wolf's general Wielandt bound $L_0\leq(D_j^2-k+1)D_j^2\leq(D^2+1)^2$
    \cite{Sanz2010Quantum,Wolf2012Quantum} and $b\leq D$ turn
    (\ref{eq:rmp_w_state_ti_bound_conditional}) into
    $N<2\max\bigl((D^2+1)^2,\,3(D-1)((D^2+1)^2+1)\bigr)$, that is,
    $D=\Omega(N^{1/5})$.\footnote{Formal declaration:
    \leanid{MPSTensor.PGVWC07CanonicalFormData.lt_of_mpv_eq_smul_wIndicator_of_isPrimitive}.}
\end{itemize}

\section{Local fix: the single-block threshold}

For $b=1$ the printed inequality $N/2\leq 3(b-1)(L_0+1)$ reads $N\leq0$. It
fails for $\ket{W_1}=\ket{1}$, the periodic vector of the single $1\times1$
block $A^0=0$, $A^1=1$ with weight $1$. This block satisfies the canonical
conditions and condition C1 at $L_0=1$. The proof itself needs both pieces to
contain at least $L_0$ sites, since condition C1 is applied to each piece;
for $b=1$ the direct-sum length $3(b-1)(L_0+1)$ is zero. The formal bound
(\ref{eq:rmp_w_state_ti_bound_conditional}) therefore uses the threshold
$\max(L_0,3(b-1)(L_0+1))$. For $b\geq2$ this equals the printed
$3(b-1)(L_0+1)$.

The step called trivial is false without a lower bound on $N$:
\begin{align}
    \ket{W_2}
        &=
        \tfrac12\bigl((\ket0+\ket1)^{\otimes2}-(\ket0-\ket1)^{\otimes2}\bigr)
    \label{eq:rmp_w_state_ti_bound_w_two}
\end{align}
is a sum of two translation-invariant product states, that is, of two distinct
$1\times1$ blocks. This does not affect the corollary: two blocks force
$N\geq 2\cdot3(L_0+1)\geq12$ in its range, where the argument above applies.

\section{What remains open}

The bound (\ref{eq:rmp_w_state_ti_bound_conditional}) is conditional. Two
inputs are missing for
(\ref{eq:rmp_w_state_ti_bound_asymptotic_lower_bound}) and for the source's
$D\succeq O(N^{1/3})$.
\begin{itemize}
    \item A reduction of an arbitrary $D\times D$ tensor whose periodic vector
    at one length $N$ is $\ket{W_N}$ to a canonical form whose blocks satisfy
    condition C1 at a controlled length. The source obtains this from
    Conjecture~2 and neglects peripheral eigenvalues of modulus one, which it
    removes by blocking or by taking $N$ prime
    \cite[lines~2090--2097]{PerezGarcia2007MPSReps}.

    \item The explicit Wielandt index $L_0=O(D^2)$ (Conjecture~2) or
    $L_0=O(D^2\log D)$ (Micha\l{}ek and Shitov). Only Wolf's general bound
    $(D^2-k+1)D^2$ is formalized, which yields $D=\Omega(N^{1/5})$ for
    primitive blocks.
\end{itemize}

\section{Formalized open-boundary representation}

The example file defines the site-independent W tensor, the two boundary
covectors, and its open-boundary contraction. For a bit string
$\sigma=(\sigma_1,\ldots,\sigma_N)$, define the amplitude
\begin{align}
    \mathcal A_N(\sigma)
        &=
        (0|\,A^{\sigma_1}\cdots A^{\sigma_N}\,|1).
    \label{eq:rmp_w_state_ti_bound_open_boundary_amplitude}
\end{align}
Let $\mathbf{1}_{\mathrm{wt}(\sigma)=1}$ denote the indicator that
$\sigma$ has exactly one nonzero entry. The formal result proves
\begin{align}
    \mathcal A_N(\sigma)
        &=
        \mathbf{1}_{\mathrm{wt}(\sigma)=1}.
    \label{eq:rmp_w_state_ti_bound_single_excitation_indicator}
\end{align}
Therefore the open-boundary contraction is exactly the unnormalized state
$\ket{W_N}$ from (\ref{eq:rmp_w_state_ti_bound_w_state}).\footnote{Formal
declarations:
\leanid{MPSTensor.wTensor},
\leanid{MPSTensor.openState}, and
\leanid{MPSTensor.wTensor_openState_eq_wIndicator}.}

\section{Formalized periodic statements}

Two weaker periodic statements are proved. First, closing the printed tensor
with a trace gives $\operatorname{tr}(A^{\sigma_1}\cdots A^{\sigma_N})=2$ on
the vacuum and $0$ on every configuration with an excitation, so the printed
tensor is not a periodic representation of $\ket{W_N}$ for $N\geq1$. Second,
let $B$ have bond dimension $D$ and satisfy
$\operatorname{tr}(B^{\sigma_1}\cdots B^{\sigma_k})=\ket{W_k}(\sigma)$ for
$1\leq k\leq D$. The vacuum amplitudes give
$\operatorname{tr}((B^0)^k)=0$ for these $k$, so the Newton--Girard identities
and the Cayley--Hamilton theorem give $(B^0)^D=0$. The configuration
$\ket{10\cdots0}$ then has amplitude $0$ at every length $N>D$, so $B$
represents no $\ket{W_N}$ with $N>D$. In particular no single tensor
represents $\ket{W_N}$ for all $N\geq1$, and for $D=1$ no tensor represents
$\ket{W_N}$ at any single length $N\geq2$.\footnote{Formal declarations:
\leanid{MPSTensor.mpv_wTensor},
\leanid{MPSTensor.not_isPeriodicWState_of_bondDim_one},
\leanid{MPSTensor.not_isPeriodicWState_of_lt}, and
\leanid{MPSTensor.not_forall_isPeriodicWState}.}
These statements constrain one tensor across several lengths. The source bound
(\ref{eq:rmp_w_state_ti_bound_asymptotic_lower_bound}) constrains the bond
dimension at a single length $N$; its formalized conditional core is described
above.

Neither
(\ref{eq:rmp_w_state_ti_bound_asymptotic_lower_bound}) nor
(\ref{eq:rmp_w_state_ti_bound_polynomial_consequence}) has been formalized.
The development does not assert either bound as an established formal result.

\section{Verdict}

This note records a deferred result, not a gap in a claimed formalization. The
open-boundary representation with $D=2$ is fully proved. The core of the
source proof is formalized as the conditional bound
(\ref{eq:rmp_w_state_ti_bound_conditional}), with the single-block threshold
corrected, and with primitive blocks it gives $D=\Omega(N^{1/5})$. The
review's bound (\ref{eq:rmp_w_state_ti_bound_asymptotic_lower_bound}) and the
source's $D\succeq O(N^{1/3})$ remain unformalized: they need a reduction of
an arbitrary tensor to a canonical form with condition C1 at a fixed length and
the explicit $O(D^2)$ or $O(D^2\log D)$ Wielandt index.

\bibliographystyle{alpha}
\bibliography{references}

\end{document}
